\documentclass[a4paper,14pt]{extarticle}
\usepackage{array,amsmath,fancybox,amsfonts,graphicx,amsthm,amssymb,wasysym,latexsym,slashed,anysize,subfigure,calc}
\usepackage[spanish]{babel}
\usepackage[utf8]{inputenc}
\usepackage{fancyhdr,float,hyperref}
\usepackage{lettrine}

\usepackage{enumitem}


%fractiones grandes en inline text: inventa un comando \ddfrac
\newcommand\ddfrac[2]{\frac{\displaystyle #1}{\displaystyle #2}}

%lewis y química
\usepackage{chemformula,elements}
%chemformula, provides \ch{} with the !(<below>)(<formula>) syntax and \chlewis;
%elements, provides \elconf and \writeelconf.
%ex:
%\ch{
%  !(\elconf{Li})( "\chlewis{180.}{Li}" ) +
%  !(\elconf{F})( "\chlewis{0.90:180:270:}{F}" )
%   ->
%  !(\writeelconf{2})( Li+ ) +
%  !(\writeelconf{2,2+6})( "\chlewis{0:90:180:270:}{F}" {}- )
%}


%entorno para química: texto y math mode
%\newcommand*\chem[1]{\textup{#1}}
\newcommand*\chem[1]{\ensuremath{\mathrm{#1}}}

%column
\usepackage{multicol}
%uso: \begin{multicols}{3}

%enumitem: configurar listas
\usepackage{enumitem}

%para los entornos de soluciones
\usepackage{mdframed,xcolor,color}
	%uso: 	entorno (begin,end): {mdframed}[backgroundcolor=brown!10] 

%subrayado y highlight. Tiene que estar detrás de usepackage[color]
\usepackage{soul}
	%\textst{normal} 
	\DeclareRobustCommand{\hlcyan}[1]{{\sethlcolor{cyan}\hl{#1}}} %en cyan!
		%\hl{highlighted text} %uso


%tamaños fuentes extra: 8,9,14,17,20
	\usepackage{extsizes}
	
 	%fuente librebaskerville
 	\usepackage[T1]{fontenc}
	\usepackage{librebaskerville}
	
	%fracciones $$ en texto más elaboradas
	 \usepackage{xfrac}
	 
	%permitir que una ecuación larga esté separada por un salto de página
	\allowdisplaybreaks[4]
		
	%gráficos varios 
	\usepackage{graphics}
	
	%fancy tables
	\usepackage{booktabs}
	\usepackage[small,centerlast,up]{caption}

	
	%mejora en la exposición de guarismos
	\usepackage{numprint}
	
\usepackage{everypage,draftwatermark} %Marca de agua
\SetWatermarkText{\vspace*{2cm}

\hspace*{2cm}\parbox{2\textwidth}{\centering DIST.\\B2-FÍSICA}}
\SetWatermarkScale{1.5}
\SetWatermarkLightness{0.9}

\DeclareMathOperator{\lin}{lin}
\DeclareMathOperator{\id}{\textbf{id}}
\DeclareMathOperator{\modulo}{mod}

\newcommand{\mc}[3]{\multicolumn{#1}{#2}{#3}}
\newcommand{\halmos}{
\vspace*{-2ex}

\begin{flushright}
\rule{1mm}{2.5mm} \end{flushright}
\vspace*{-3ex}

}
\renewcommand{\qed}{\halmos}


%versores: %uso: \uveci\ne\hat{\imath}_{\uvec{i}}+\uvec{j}+\uvec{k}$
\usepackage{bm}
\newcommand{\uveci}{{\bm{\hat{\textnormal{\bfseries\i}}}}}
\newcommand{\uvecj}{{\bm{\hat{\textnormal{\bfseries\j}}}}}
\DeclareRobustCommand{\uvec}[1]{{%
  \ifcsname uvec#1\endcsname
     \csname uvec#1\endcsname
   \else
    \bm{\hat{\mathbf{#1}}}%
   \fi
}}


%algún nuevo operador
\DeclareMathOperator{\cotan}{cotan}


%para textbf en entorno enumerate
\renewcommand{\labelenumi}{%
 \textbf{\theenumi}.
}

%a4
\marginsize{.5cm}{1cm}{1cm}{1cm}%requiere el paquete anysize,izq,dcha,arriba,abajo

%a5
%\marginsize{1cm}{1cm}{1cm}{.5cm}%requiere el paquete anysize,izq,dcha,arriba,abajo

\newtheorem{teo}{{\sc \textbf{Teorema}}}
\newtheorem{defi}{{\sc \textbf{Definición}}}
\newtheorem{coro}{{\sc \textbf{Corolario}}}


%opening
\title{\textbf{	}}
\date{}
\pagestyle{fancy}
\fancyhead[LO,RE]{{\small DOSSIER DE PROBLEMAS\hspace{1ex} \textbf{2Ev}}}
\fancyhead[RO,LE]{{\small B2-DIST\hspace{2ex} FÍSICA}}
%\fancyhead[CO,CE]{Ricardo Torres\\http://tinyurl.com/iesCMGfyq}



\begin{document}
\pagenumbering{gobble}

%gas ideal y mezclas
%disoluciones
%estequiometría: ajuste, reactivo limitante, reactivos en disolución e impuros


\begin{enumerate}
  \setlength\itemsep{1.5em}
  
   \item (2,5 puntos)	Un sistema está formado por dos lentes convergentes \textit{iguales} separadas $70\;cm$, con $20\;cm$ de distancia focal. Un objeto se coloca a $40\;cm$ de la primera lente. 
		\begin{itemize}
			\item[a)] Calcula la \textbf{posición de la imagen} a través del sistema completo y determina el \textbf{aumento lateral total}.
			\item[b)]  Realiza el \textbf{trazado de rayos} correspondiente.			%\vspace*{.1cm}
			
			 \qed
		\end{itemize}
    
    
    \item  
		\begin{itemize}
			\item[a)] (1,25 puntos) Aproxima el valor del \textbf{radio de curvatura} de las superficies de una lente \textit{simétrica} biconvexa de $4$ dioptrías sabiendo que el índice de refracción del vidrio que forma la lente es $n\simeq 1.5$.
			%25cm
			\item[b)] (1,25 puntos) Una persona \textit{miope} utiliza gafas con lentes divergentes de $2.5$ dioptrías  para ver correctamente objetos muy alejados.  Estima la \textbf{distancia máxima} a la que puede enfocar nítidamente sin gafas.
			%40 cm
			\qed
		\end{itemize}				

    
    \item (2,5 puntos)  Puliendo  por  frotamiento  una  de  las  caras  de  un  cubito  de  hielo,  puede construirse una lente convergente \textit{plano-convexa}.
		\begin{itemize}
			\item[a)] Estima el  \textbf{radio  de  curvatura}  que debería tener la  cara  pulida para que pudiera ser utilizada, en una urgencia, por una persona \textit{présbita} que necesita unas gafas de $5$ dioptrías para poder leer.
			%r=6 cm
			\item[b)]  Digan lo que digan respecto a la mala relación del agua y el fuego, resulta que el \textit{cubito-lente}  también puede emplearse   para   encender una hoguera. Si quisieras prender fuego a un papel a partir de los rayos del Sol, \textit{razona} \textbf{en qué posición y a qué distancia} lo colocarías del \textit{cubito-lente} calculado en el apartado \textit{a)}.
			%f'=20 cm
			\end{itemize}
				
					\textbf{Dato}: $n_{\text{hielo}}\simeq 1.3$		
			 \qed
    
    \item (2,5 puntos) Desde un punto a una distancia $s_0\!=\!-5f'$ a la izquierda del plano de incidencia de una lente delgada convergente de distancia focal imagen $f'$ un cochecito de juguete se acerca con una \textit{velocidad constante} $v=1\;cm\,s^{-1}$. 
			\begin{itemize}
				\item[a)] Calcula la \textbf{posición de la imagen} del coche a través del sistema óptico en cualquier instante anterior a la colisión con la lente.
				\item[b)] ¿Para qué \textbf{valores de $\mathbf{s}$} la imagen es \textit{aumentada}? 
			\end{itemize}
			
					\textbf{Dato:} $f'=2\;cm$
				
			\qed

    %en parcial P2 de 2021 (P2_BC/p2_sol)
    
\end{enumerate}

\end{document}

